% !TEX root = ../main.tex

\chapter{总结与展望}

\section{研究工作总结}

本文围绕“基于 Flink 的动态 Kafka Topic 发现与订阅机制”这一主题，完成了从问题提出、相关工作分析，到系统设计、实现与验证的完整过程。
与仅停留在概念层面的方案不同，本文工作已经落到可运行的连接器原型：
核心实现覆盖元数据发现、运行时路由、写入通道管理、路由级状态快照和提交协调等能力；
实验载体则用于驱动动态发现、显式路由更新与静态基线对比。

从实现路径看，本文并未完全推翻现有 KafkaSink，而是采取“动态外壳 + 静态内核”的思路：
外层负责路由解析、元数据刷新与多写入通道协调，内层每个物理路由仍复用标准 \id{KafkaSink} 的写入、状态与提交逻辑。
这种设计既满足了本文关于“逻辑路由与物理路由分离”的设计目标，也体现了“基于正则表达式实现动态发现、
并与 Flink Checkpoint 机制结合”的问题界定。当前实现已经能够在不重启作业的情况下接收路由更新事件、
自动解析匹配 Topic 并完成多目标写入，形成了一套完整可运行的动态 Sink 方案。

进一步而言，本文的价值并不只在于“实现了一个能运行的 demo”，更在于把问题拆分为可操作的几个关键环节：
一是以元数据服务和正则匹配为基础完成动态目标发现；二是以逻辑路由与物理 Topic 解耦完成运行期路由更新；
三是以路由为单位组织写入通道、状态和提交路径，使动态行为尽可能兼容现有 KafkaSink 的实现框架。
这种拆分方式使课题既具备工程可落地性，也具备较清晰的方法论结构。

\section{存在的不足}

尽管当前 connector 已具备原型可运行性，但距离更成熟的工程化方案仍有若干差距。首先，
当前元数据发现依赖 Kafka 管理面周期性轮询，在 Topic 规模扩大后会带来额外管理面开销，并且 discovery interval 决定了发现时延下界。
其次，当前实现能够新增路由对应的写入通道，但在长时间运行或频繁变更场景下，退役通道的安全回收仍需要更细致的策略。
再次，当前验证载体主要覆盖至少一次语义，虽然动态 Sink 设计上支持 exactly-once，但尚未结合复杂故障场景完成系统化验证。

此外，当前自动发现路径默认会把普通记录广播写入所有活动路由，这在某些“多目标复制”场景中是合理的，
但在更一般的业务分流场景下可能造成不必要的写放大。日志、监控指标、异常治理与配置版本控制等工程特性，
也仍需进一步加强。这些不足说明当前原型在元数据访问开销、实时性、一致性和工程化能力方面仍有完善空间。

在本文最新实现中，逻辑路由的删除语义已经补充完成，非事务路径下的退役写入通道清理也已落地，
这使路由生命周期管理更加完整；但 exactly-once 路径上的安全回收、事务边界验证与
更精细的资源治理仍未完全解决。需要强调的是，这些不足并不否定原型本身的研究意义，而是说明当前实现更适合作为“动态 Sink 机制的验证性实现”。
也就是说，本文已经证明了动态发现、动态路由与路由级状态封装在 Flink Kafka Sink 场景下具有可行性，
但距离一个大规模、长期稳定运行的生产级方案，还需要在资源治理、边界条件和异常恢复方面继续完善。

\section{未来工作展望}

未来的改进可从四个方向推进。第一，优化元数据访问机制，在保持动态发现能力的同时降低轮询成本，
例如引入缓存、差分刷新或更细粒度的发现策略，以缩短 Topic 变化到可写之间的时延。第二，
完善写入通道生命周期管理，在路由失活后安全回收对应通道，并研究更合理的路由级资源上限控制策略。
第三，围绕 exactly-once 路径设计更系统的故障恢复实验，验证路由级状态恢复、事务前缀管理、
提交重试与幂等性边界。第四，进一步强化工程化能力，包括指标埋点、日志可观测性、异常分类处理、
配置热更新版本管理以及更丰富但不改变 connector 核心语义的验证载体。

在论文工作层面，后续还需要结合完整实验数据补充第 5 章中的定量分析，例如 discovery interval 对发现延迟的影响、
路由数量增长对吞吐与资源占用的影响，以及动态方案与静态 \id{KafkaSink} 的对比结果。通过这些补充，
可使本文从“原型验证与可行性说明”进一步提升为“具有更强实证支撑的毕业设计成果”。

总体而言，本文已经完成了毕业设计从“选题论证”到“原型落地”的关键阶段。后续只要继续围绕实验数据补充、
实现细节打磨和论文语言精修推进，就可以使整篇论文在工程深度、论证完整性和学术表达上进一步趋于成熟。
